% Template for a single-paper introduction.
% Compile from slide/template/ (PDF: slide/out/template/ , other files: build/):
%   ../lualatex-dated.sh -interaction=nonstopmode -output-directory=out template_en.tex
% Colors, title page, and the takeaway bar live in labslide.sty.
% Frame titles use "Topic | detail". End with \takeaway then \source.
% Replace every bracketed placeholder.
\documentclass[aspectratio=169,10pt,t]{beamer}
\usepackage[en]{labslide}

\kicker{Paper introduction}
\title[Paper introduction template]{Paper Introduction Slides}
\subtitle{A template: one paper, from problem to method to results}
\author{[Authors] ([Affiliation])}
\venue{[Venue] [Year]}
\date{[Date]}

\begin{document}

\begin{frame}[plain]
  \labtitlepage{%
    \agendacard[blue]{What is hard}{Name the scene in one sentence.}%
    \hfill
    \agendacard[orange]{What changes}{Inputs, training, and deployment.}%
    \hfill
    \agendacard[green]{What holds}{The baseline, and what remains.}%
  }
\end{frame}

\begin{frame}{Problem | The situation this paper takes on}
  State the robot, the environment, and the constraint in one sentence.

  \begin{itemize}
    \item Constraint 1: time, sensing, compute, or safety.
    \item Constraint 2: the information available at training differs from deployment.
    \item Constraint 3: the cost of a failure while collecting data.
  \end{itemize}

  \takeaway{Say whether the paper trains the gait, the adaptation, or both.}
  \source{[Authors], [year], Sec.~1}
\end{frame}

\begin{frame}{Related work | Where the approaches differ}
  \begin{columns}[T,onlytextwidth]
    \begin{column}{0.32\textwidth}
      \infocard[blue]{[Line A]}{What is designed by hand. Where it holds, and where it breaks.}
    \end{column}
    \hfill
    \begin{column}{0.32\textwidth}
      \infocard[orange]{[Line B]}{What is randomized for robustness. The performance that is given up.}
    \end{column}
    \hfill
    \begin{column}{0.32\textwidth}
      \infocard[green]{[Line C]}{What is collected at deployment. How many trials that takes.}
    \end{column}
  \end{columns}

  \takeaway{State which missing piece of which line this paper fills.}
  \source{[Authors], [year], Sec.~2}
\end{frame}

\begin{frame}{Comparison | Same goal, different place to intervene}
  \begin{columns}[T,onlytextwidth]
    \begin{column}{0.48\textwidth}
      \infocard[orange]{[Method A]}{%
        Cue: [input].\\[0.25em]
        Aim: [quantity improved].\\[0.25em]
        Left open: [next question].}
    \end{column}
    \hfill
    \begin{column}{0.48\textwidth}
      \infocard[blue]{[Method B]}{%
        Cue: [input].\\[0.25em]
        Aim: [quantity improved].\\[0.25em]
        Left open: [next question].}
    \end{column}
  \end{columns}

  \takeaway{Describe the difference as a difference of information, then of outcome.}
  \source{[Paper A] / [Paper B]}
\end{frame}

\begin{frame}{Method | Split the inputs used in training and at runtime}
  \centering
  \vspace{1.2em}
  {\footnotesize\color{muted} During training}\\[0.45em]
  \flowbox{Environment\\{[$e_t$]}}
  \flowarrow
  \flowbox{Encoder\\{[$\mu$]}}
  \flowarrow
  \flowbox{Latent\\{[$z_t$]}}
  \flowarrow
  \flowbox{Policy\\{[$\pi$]}}
  \flowarrow
  \flowbox{Action\\{[$a_t$]}}

  \vspace{1.5em}
  {\footnotesize\color{muted} At deployment}\\[0.45em]
  \flowbox{Recent\\history}
  \flowarrow
  \flowbox{Estimator\\{[$\phi$]}}
  \flowarrow
  \flowbox{Estimate\\{[$\hat z_t$]}}
  \flowarrow
  \flowbox{Same policy\\{[$\pi$]}}

  \takeaway{State which input exists only in training, and what replaces it at runtime.}
  \source{[Authors], [year], Fig.~X / Sec.~3}
\end{frame}

\begin{frame}{Definitions | Give each symbol a dimension and a role}
  \[
    z_t = \mu(e_t), \qquad a_t = \pi(x_t, a_{t-1}, z_t)
  \]

  \begin{itemize}
    \item $x_t$: [contents and dimension of the state].
    \item $a_t$: [contents of the action], and how it becomes torque or position.
    \item $e_t$: [quantities available only in training]. $z_t$ is their summary.
  \end{itemize}

  \takeaway{Say whether the estimate is a physical parameter or the summary the policy uses.}
  \source{[Authors], [year], Sec.~3.1}
\end{frame}

\begin{frame}{Results | Hold the comparison to the adaptation itself}
  \vspace{0.4em}
  \centering
  {\footnotesize
  \begin{tabular}{lccc}
    \toprule
    Method & Success (\%) & Time to fall & Distance \\
    \midrule
    {[Baseline A]} & -- & -- & -- \\
    {[Baseline B]} & -- & -- & -- \\
    Proposed & -- & -- & -- \\
    \midrule
    {[Upper bound]} & -- & -- & -- \\
    \bottomrule
  \end{tabular}}

  \vspace{0.8em}
  \raggedright
  State the test condition in one line: [does the environment change mid-episode, and are real-world samples required].

  \takeaway{Read each number against the upper bound and against the extra data it needs.}
  \source{[Authors], [year], Table~X}
\end{frame}

\begin{frame}{Summary | Name the question that is still open}
  \begin{itemize}
    \item The policy is trained with [privileged inputs], and those inputs are recovered at runtime from [history or sensors].
    \item Left unused: reference motions, extra real-world learning, and simulator calibration.
    \item Evidence: [simulation comparison] and [real scenes].
    \item Limit: [when it fails], and the next sensor that would address it.
  \end{itemize}

  \takeaway{Leave one sentence: which module adapts, and from which information.}
  \source{[Authors], [year], conclusion}
\end{frame}

\begin{frame}{References}
  {\small
  \begin{enumerate}
    \item {[Authors]. [Title]. [Venue], [year].}\\
      \url{https://arxiv.org/abs/0000.00000}
    \item {[First author] et al. [Title]. [Venue], [year].}
  \end{enumerate}}
  \vfill
  {\scriptsize\color{muted}
  Match the section numbers used in the talk. First author plus ``et al.'' is enough.}
\end{frame}

\end{document}
